按 2022 年公报口径，问题二以江豚 1249 头、长江鲟 438 尾作为初值与口径锚定。基线修复情景下，江豚和长江鲟的归一化末值分别为 1.133 和 1.140。仅提高通道阻隔且不叠加污染时，两者的末值分别降至 0.901 和 0.856，较基线下降 20.5\% 和 24.9\%。将人工放流项置为 \(R_S=0\) 后，江豚末值为 1.145，较基线变化 +1.0\%；长江鲟末值为 0.519，较基线下降 54.5\%。污染情景下，两者的末值分别为 0.623 和 0.706，较基线下降 45.0\% 和 38.0\%。

\(\delta_B\) 表示通道阻隔和栖息地碎片化的综合损失。仅提高通道阻隔时，江豚和长江鲟的末值均低于基线；污染情景则用于比较污染叠加阻隔对基线收益的侵蚀。人工放流项置零后，江豚末值未下降，长江鲟末值降幅较大，这一对照用于拆分自然恢复与人工放流补偿效应。上述认识限于当前参数组与 2022 公报归一化口径下的内部投影，不构成独立外部验证；\(M\) 和 \(V\) 是传导和竞争状态，不是本问的直接观测目标。
